\begin{vidu}
Một con lắc lò xo treo thẳng đứng gồm lò xo có độ cứng $k=100\ N/m$ và chiều dài tự nhiên $\ell_0=40\ cm$. Đầu trên của lò xo gắn cố định vào tường còn đầu dưới gắn vào vật có khối lượng $m=500\ g$. Kích thích cho vật dao động điều hoà với biên độ $A=10\ cm$. Lấy $g=10\ m/s^2$.
\begin{itemize}
\item[1.] Độ biến dạng của lò xo khi vật ở vị trí cân bằng là ........ cm
\item[2.] Chiều dài của lò xo khi vật ở vị trí cân bằng là ........ cm
\item[3.] Chiều dài lớn nhất của lò xo là ........ cm
\item[4.] Chiều dài nhỏ nhất của lò xo là ........ cm
\item[5.] Khi vật đạt độ cao cực đại thì lò xo nén ........ cm
\item[6.] Khi vật ở vị trí thấp nhất thì lò do dãn ........ cm
\item[7.] Độ lớn lực đàn hồi cực đại trong quá trình vật dao động là ........
\item[8.] Độ lớn lực đàn hồi cực tiểu trong quá trình vật dao động là ........
\item[9.] Khoảng thời gian ngắn nhất giữa hai lần liên tiếp lò xo có chiều dài $50\ cm$ là ........
\item[10.] Khoảng thời gian lò xo bị nén trong một chu kì là ........
\item[11.] Khoảng thời gian lò xo bị dãn trong một chu kì là ........
\end{itemize}
\loigiai{
\begin{itemize}
\item[1.] Độ biến dạng của lò xo khi vật ở vị trí cân bằng là 5 cm
\item[2.] Chiều dài của lò xo khi vật ở vị trí cân bằng là 45 cm
\item[3.] Chiều dài lớn nhất của lò xo là 55 cm
\item[4.] Chiều dài nhỏ nhất của lò xo là 35 cm
\item[5.] Khi vật đạt độ cao cực đại thì lò xo nén 5 cm
\item[6.] Khi vật ở vị trí thấp nhất thì lò do dãn 15 cm
\item[7.] Độ lớn lực đàn hồi cực đại trong quá trình vật dao động là 15 N
\item[8.] Độ lớn lực đàn hồi cực tiểu trong quá trình vật dao động là 5 N
\item[9.] Khoảng thời gian ngắn nhất giữa hai lần liên tiếp lò xo có chiều dài $50\ cm$ là $\dfrac{\sqrt{5}}{15}\ N\approx 0,15\ s$
\item[10.] Khoảng thời gian lò xo bị nén trong một chu kì là $\dfrac{\sqrt{5}}{15}\ N\approx 0,15\ s$
\item[11.] Khoảng thời gian lò xo bị dãn trong một chu kì là $\dfrac{2\sqrt{5}}{15}\ N\approx 0,3\ s$
\end{itemize}
}
\end{vidu}
%\dotlineEX{13}
\begin{vidu}%[3P1-6.2B][Loai1]  
Một con lắc lò xo treo thẳng đứng, độ cứng $k=100$ N/m, vật nặng khối lượng $m=500\ g$.  Lấy $g = 10\ m/s^2$. Khi vật cân bằng lò xo dãn
\choice
{\True 5 cm}
{2 cm}
{4 cm}
{2,5 cm}
\loigiai{
Độ dãn của lò xo tại vị trí cân bằng là: $\Delta \ell_0=\dfrac{mg}{k}=\dfrac{0,5.10}{100}=5\ cm$.
}
%\haidong{6}
\end{vidu} %\dotlineEX{2}
%\loai{Tìm chu kì}
\begin{vidu}%[3P1-6.2B][Loai1]
Một con lắc lò xo treo thẳng đứng gồm lò xo có chiều dài tự nhiên là 20 cm. Tại vị trí cân bằng lò xo dài 24 cm. Lấy $g = 10\ m/s^2$. Chu kì dao động riêng của con lắc lò xo này là
\choice
{\True 0,397 s}
{1 s}
{2 s}
{1,414 s}
\loigiai{
Ta có: $\Delta \ell_0  = \ell_{cb} - \ell_0 = 4\ cm \Rightarrow  \omega =\sqrt{\dfrac{g}{\Delta \ell_0}}=   5\sqrt{10}\ rad/s \Rightarrow T= 0,397\ s$.
}
%\haidong{5}
\end{vidu} %\dotlineEX{4}
%\loai{Tìm chiều dài lò xo ở vị trí cân bằng}

%\loai{Tìm chiều dài tự nhiên của lò xo}
\begin{vidu}%[3P1-6.2B][Loai1]
Một con lắc lò xo treo thẳng đứng dao động điều hòa với chu kì 0,4 s. Khi vật ở vị trí cân bằng, lò xo dài 44 cm. Lấy $g = \pi^2 \ m/s^2$. Chiều dài tự nhiên của lò xo là
\choice
{36 cm}
{\True 40 cm}
{42 cm}
{38 cm}
\loigiai{
Ta có: $T = 0,4\ s \Rightarrow  \Delta \ell_0 =\dfrac{g}{\omega^2}=  4\ cm \Rightarrow \ell_0 = \ell_{cb} - \Delta \ell_0  = 40\ cm$.
}
%\haidong{4}
\end{vidu} %\dotlineEX{4}
%\loai{Tìm chiều dài max, min của lò xo}
\begin{vidu}%[3P1-6.2K][Loai1]
Con lắc lò xo treo thẳng đứng, dao động điều hòa với phương trình $x=4\cos  \left( 5\pi t+\dfrac{\pi}{6} \right)$ cm. Lấy $g = 10\ m/s^2$ và $\pi^2 = 10$. Chiều dài tự nhiên của lò xo là 50 cm. Chiều dài lớn nhất và nhỏ nhất của lò xo trong quá trình dao động của vật lần lượt là
\choice
{60 cm và 52 cm}
{60 cm và 54 cm}
{\True 58 cm và 50 cm}
{56 cm và 50 cm}
\loigiai{
\begin{itemize}
\item Ta có: $ \Delta \ell_0=\dfrac{g}{\omega^2}= \ 4\ cm \Rightarrow \ell_{cb} = \ell_0 + \Delta \ell_0  = 54\ cm \Rightarrow \ell_{\max} = \ell_{cb} + A = 58\ cm$ và $\ell_{\min} = \ell_{cb} - A = 50\ cm$.
\item 
\item 
\end{itemize}
}
%\haidong{6}
\end{vidu} %\dotlineEX{5}

%\loai{Biết chiều dài max, min tìm A, $\Delta \ell_0$ và năng lượng}
\begin{vidu}%[3P1-6.2K][Loai1]  
Một con lắc lò xo đặt nằm ngang, lò xo khối lượng không đáng kể có độ cứng $100\ N/m$, dao động điều hòa. Trong quá trình dao động chiều dài của lò xo biến thiên từ 20 cm đến 32 cm. Cơ năng của vật là
\choice
{3 J}
{\True 0,18 J}
{1,5 J}
{0,36 J}
\loigiai{
\begin{itemize}
\item Biên độ dao động của con lắc $A=\dfrac{\ell_{\max }-\ell_{\min}}{2}=\dfrac{32-20}{2}=6\ cm$.
\item Cơ năng của vật $W=\dfrac{1}{2}{kA}^2=\dfrac{1}{2}.100.0,06^2=0,18\ J$.
\end{itemize}
}
%\haidong{6}
\end{vidu} %\dotlineEX{2}
%\loai{Tìm chiều dài lò xo ở li độ và thời điểm cho trước}
\begin{vidu}%[3P1-6.2K][Loai1]
Con lắc lò xo treo thẳng đứng dao động theo phương trình $x=A\cos  (\omega t+\varphi )\ cm.$ Chọn chiều dương hướng xuống. Khi con lắc dao động thì lò xo có $\ell_{\max} = 1$ m và $\ell_{\min} = 0,8$ m. Khi pha dao động là $\dfrac{2\pi}{3}$ thì chiều dài lò xo là
\choice
{\True 85 cm}
{90 cm}
{87,5 cm}
{92,5 cm}
\loigiai{
\begin{itemize}
\item Ta có: $\ell_{cb}=\dfrac{\ell_{\max}+\ell_{\min}}{2}=90 \ cm;\   A=\dfrac{\ell_{\max}-\ell_{\min}}{2}=10 \ cm$.
\item Khi pha dao động $\phi =\dfrac{2\pi}{3}  \Rightarrow  x=-\dfrac{A}{2}(-)=-5\ cm\quad (-) $, chiều dương hướng xuống $\Rightarrow$ vật đang ở trên vị trí cân bằng 5 cm
$\Rightarrow \ell  = \ell_{cb} - 5= 85\ cm$.
\end{itemize}
}
%\haidong{6}
\end{vidu} %\dotlineEX{3}
%\loai{Tìm tốc độ cực đại}
\begin{vidu}%[3P1-6.2K][Loai1]
Một con lắc lò xo treo thẳng đứng. Kéo vật xuống dưới vị trí cân bằng 3 cm rồi truyền cho nó tốc độ 40 cm/s thì nó dao động điều hòa theo phương thẳng đứng trùng với trục của lò xo. Khi vật đạt độ cao cực đại, lò xo dãn 5 cm. Lấy gia tốc trọng trường $g = 10\ m/s^2$. Tốc độ cực đại của vật dao động là
\choice
{1,15 m/s}
{\True 0,5 m/s}
{10 cm/s}
{2,5 cm/s}
\loigiai{
%\begin{itemize}
%\item Khi vật ở độ cao cực đại (biên trên) lò xo dãn 5 cm $\Rightarrow$ vật ở dưới vị trí lò xo tự nhiên 5 cm $\Rightarrow \Delta \ell_0 = A + 5\ cm$.
%\item Khi truyền tốc độ vật có $|x| = 0,03\ m$ và $|v| = 0,4\ m/s\Rightarrow  A=\sqrt{x^2+\dfrac{v^2}{\omega^2}}\Rightarrow A^2=x^2+\dfrac{v^2\Delta \ell_0}{g}\\
%\Leftrightarrow \left( 0,05-\Delta \ell_0 \right)^2=0,03^2+\dfrac{0,4^2\Delta \ell_0}{10} \Rightarrow \Delta \ell_0 = 1,6\ cm$ và  $\Delta \ell_0 = 10\ cm$.
%\item Khi vật cao nhất lò xo dãn 5 cm $\Rightarrow \Delta \ell_0 = 10\ cm$ và $A = 5\ cm$.
%$\omega =\sqrt{\dfrac{g}{\Delta \ell_0}}  = 10\ rad/s \Rightarrow v_{\max} = \omega A = 50\ cm/s = 0,5\ m/s$.
%\end{itemize}
}
%\haidong{7}
\end{vidu} %\dotlineEX{4}
%\loai{Viết phương trình dao động}

%\loai{Tìm tốc độ tại vị trí bất kì}
\begin{vidu}%[3P1-6.2K][Loai1]
Một con lắc lò xo treo thẳng đứng gồm lò xo có chiều dài tự nhiên là 40 cm và vật nhỏ đang dao động điều hòa. Lò xo có chiều dài biến thiên từ 36 cm đến 52 cm trong quá trình dao động. Lấy $g = \pi^2 \ m/s^2$. Tốc độ của vật khi lò xo dài 40 cm là
\choice
{\True $20\pi \sqrt{3}\ cm/s$}
{$40\pi\ cm/s$}
{$20\pi \sqrt{2}\ cm/s$}
{$20\pi\ cm/s$}
\loigiai{
\begin{itemize}
\item Ta có: $\ell_{cb}=\dfrac{\ell_{\max}+\ell_{\min}}{2}=44 \ cm;\  A=\dfrac{\ell_{\max}-\ell_{\min}}{2}=8 \ cm$.
\item Mặt khác: $ \Delta \ell_0=\ell_{cb}-\ell_0=4 \ cm \Rightarrow \omega  = 5\pi\ rad/s$.
\item Khi lò xo dài 40 cm $\Rightarrow$ vật cách vị trí cân bằng $|x| = |40 - \ell_{cb}| = 4\ cm$\\ $\Rightarrow \left| v \right|=\omega \sqrt{A^2-x^2} = 20\pi \sqrt{3}\ cm/s$.

\end{itemize}
}
%\haidong{7}
\end{vidu} %\dotlineEX{5}















%\loai{Tính lực đàn hồi cực đại,cực tiểu}
\begin{vidu}%[3P1-7.2K][Loai1]
Con lắc lò xo treo thẳng đứng gồm lò xo có độ cứng 40 N/m và vật nhỏ dao động điều hòa. Tại vị trí cân bằng lò xo dãn 10 cm. Trong quá trình dao động, chiều dài lò xo biến thiên từ 100 cm đến 110 cm. Độ lớn lực đàn hồi cực đại trong quá trình vật dao động là
\choice
{200 N}
{600 N}
{\True 6 N}
{60 N}
\loigiai{
\begin{itemize}
\item $ A=\dfrac{{{\ell}_{\max}}-{{\ell}_{\min}}}{2}=5 $ cm $\Rightarrow$ $\left|F\right|_{\max} = k(\Delta \ell + A) = 6$ N (biên dưới). 
\end{itemize}
}
%\haidong{6}
\end{vidu} %\dotlineEX{4}
%\loai{Tính tỉ số lực đàn hồi}
\begin{vidu}%[3P1-7.2K][Loai1]%Bài 15 dạng 2
Con lắc lò xo treo thẳng đứng gồm lò xo có khối lượng không đáng kể gắn với viên bi. Viên bi đang ở vị trí cân bằng thì được kéo xuống dưới theo phương thẳng đứng một đoạn 3 cm rồi thả nhẹ cho nó dao động. Hòn bi thực hiện 50 dao động mất 20 s. Cho $g = \pi^2 = 10\ m/s^2$. Tỉ số độ lớn lực đàn hồi cực đại và lực đàn hồi cực tiểu của lò xo khi dao động là
\choice
{5}
{4}
{\True 7}
{3}
\loigiai{
\begin{itemize}
\item Ta có: $T = 0,4\ s \Rightarrow \Delta \ell_0 = 0,04\ m \Rightarrow \dfrac{{{\left| F \right|}_{\max}}}{{{\left| F \right|}_{\min}}}=\dfrac{\Delta \ell_0+A}{\Delta \ell_0-A}=7$.

\end{itemize}
}
%\haidong{9}
\end{vidu} %\dotlineEX{6}
%\loai{Từ tỉ số lực đàn hồi tìm các đại lượng}
\begin{vidu}%[3P1-7.2K][Loai1]%Bài 15 dạng 2
Một con lắc lò xo dao động điều hòa theo thẳng đứng với biên độ 10 cm. Tỉ số giữa độ lớn lực đàn hồi cực đại và cực tiểu tác dụng vào điểm treo trong quá trình dao động là $\dfrac{7}{3}$. Lấy $g = \pi^2 = 10\ m/s^2$. Tần số dao động là
\choice
{\True 1 Hz}
{0,5 Hz}
{0,25 Hz}
{0,75 Hz}
\loigiai{
\begin{itemize}
\item $\dfrac{{{\left| F \right|}_{\max}}}{{{\left| F \right|}_{\min}}}=\dfrac{\Delta \ell+A}{\Delta \ell_0-A}=\dfrac{7}{3} \Rightarrow \Delta \ell_0 = 2,5A = 25$ cm $\Rightarrow \omega  = 2\pi$   rad/s $\Rightarrow$ f = 1 Hz.

\end{itemize}
}
%\haidong{8}
\end{vidu} %\dotlineEX{6}

%\loai{Từ lực đàn hồi tại vị trí bất kì tìm các đại lượng}
\begin{vidu}%[3P1-7.2G][Loai1] 
Một con lắc lò xo có chiều dài tự nhiên là 20 cm gắn vào một vật nặng có khối lượng $m = 200$ g dao động điều hòa theo phương thẳng đứng. Lấy $g = 10\ m/s^2$. Khi lò xo có chiều dài 18 cm thì vận tốc của vật bằng 0 và lực đàn hồi của lò xo có độ lớn 2 N. Năng lượng dao động của vật là
\choice
{0,1 J}
{0,04 J}
{\True 0,08 J}
{0,02 J}
\loigiai{
\begin{itemize}
\item Chiều dài tự nhiên của lò xo là $\ell_0= 20\ cm$.
\item Khi lò xo có chiều dài 18 cm thì vận tốc của vật bằng 0 $\Rightarrow $ khi đó vật đang ở vị trí biên trên:
$\ell=\ell_0+\Delta \ell_0-A\Rightarrow A-\Delta \ell_0=20-18=2\ cm$.
\item Lực đàn hồi lúc đó $F=k\left|{\Delta \ell_0-A}\right|=k.0,02=2\Rightarrow k=100\ N/m$.
\item Độ biến dạng của lò xo ở vị trí cân bằng là $\Delta \ell_0=\dfrac{mg}{k}=\dfrac{{0,2.10}}{{100}}=0,02\ m$.
\item Biên độ của con lắc là $A-\Delta \ell_0=2\Rightarrow A=4\ cm$.
\item Cơ năng của con lắc là $W=\dfrac{1}{2}{kA}^2=\dfrac{1}{2}.100.0,04^2=0,08\ J$.
\end{itemize}
}
%\haidong{6}
\end{vidu} %\dotlineEX{6}
%\loai{Lực đàn hồi tổng hợp}














%nên làm ví dụ mẫu tính toán mấy cái dễ dùng làm cùng lí thuyết luôn
%\loai{Tìm thời gian, thời điểm theo độ biến dạng}

%\loai{Tìm thời gian, thời điểm theo chiều dài}
\begin{vidu}%[3P1-10.1K][Loai1]
Một con lắc lò xo treo thẳng đứng gồm lò xo có chiều dài tự nhiên là 40 cm và vật nhỏ đang dao động điều hòa. Trong quá trình dao động, lò xo có chiều dài biến thiên từ 36 cm đến 52 cm. Lấy $g = \pi^2 \ m/s^2$. Khoảng thời gian ngắn nhất giữa hai lần liên tiếp lò xo dài 48 cm là
\choice
{$\dfrac{1}{15}$ s}
{\True $\dfrac{2}{15}$ s}
{0,2 s}
{0,1 s}
\loigiai{
\begin{itemize}
\item $\ell_{cb}=\dfrac{\ell_{\max}+\ell_{\min}}{2}=44\ cm;\ A=\dfrac{\ell_{\max}-\ell_{\min}}{2}=8\ cm$.
\item $ \Delta \ell_0=\ell_{cb}-\ell_0=4\ cm \Rightarrow\omega  = 5\pi\ rad/s$.
\item Khi lò xo dài 48 cm $\Rightarrow$ vật nằm bên dưới vị trí bân bằng đoạn $|x| = |48 - \ell_{cb}  | = 4\ cm=\dfrac{A}{2}$.
\item Khoảng thời gian ngắn nhất giữa hai lần có chiều dài 48 cm là $ \dfrac{T}{6}+\dfrac{T}{6}=\dfrac{T}{3}=\dfrac{2}{15}\ s $ (từ vị trí này vật xuống biên dưới rồi trở lại vị trí này).

\end{itemize}
}
%\haidong{11}
\end{vidu} %\dotlineEX{6}
%\loai{Tìm khoảng thời gian giữa hai lực đàn hồi cho trước}
\begin{vidu}%[3P1-10.1K][Loai1]
Một con lắc lò xo treo thẳng đứng dao động với biên độ $5\sqrt{2}$ cm và chu kì T. Khoảng thời gian ngắn nhất từ lúc lực đàn hồi có độ lớn cực đại đến lúc lực đàn hồi có độ lớn cực tiểu là 0,375T. Lấy $g = 10\ m/s^2$. Tốc độ của vật nặng khi lò xo bị nén 1 cm là
\choice
{87,6 cm/s}
{83,1 cm/s}
{57,3 cm/s}
{\True 52,92 cm/s}
\loigiai{
\begin{itemize}
\item Nếu $A < \Delta \ell_0$ thì khoảng thời gian ngắn nhất đó phải là $ \dfrac{T}{2} $ (vật đi từ biên dưới lên biên trên) $\Rightarrow$ loại.
\item Do đó, $A > \Delta \ell_0$. Theo trục phân bố thời gian $\Rightarrow \Delta \ell_0=\dfrac{A\sqrt{2}}{2}  = 5\ cm \Rightarrow \omega  = 10\sqrt{2} $ rad/s.
\item Khi lò xo bị nén 1 cm, tức trên vị trí lò xo tự nhiên 1 cm thì vật có li độ $|x| = 6\ cm\Rightarrow |v| =  20\sqrt{7}\ cm/s$.

\end{itemize}
}
%\haidong{10}
\end{vidu} %\dotlineEX{5}
















%\loai{Khoảng thời gian lò xo bị nén trong một chu kì}
\begin{vidu}%[3P1-10.2K][Loai1] 
Một con lắc lò xo $m = 200$ g, $k = 80$ N/m treo thẳng đứng. Đưa vật dọc theo trục của lò xo tới vị trí lò xo nén 1,5 cm rồi buông nhẹ cho vật dao động. Cho $g=10\ m/s^2$ và bỏ qua mọi ma sát. Quãng thời gian lò xo bị nén trong một chu kì dao động là
\choice
{0,1936 s}
{0,1250 s}
{0,0254 s}
{\True 0,0896 s}
\loigiai{
\begin{itemize}
\immini{\item Ta có: $\omega = \sqrt{\dfrac{k}{m}} = \sqrt{\dfrac{80}{0,2}} = 20\ rad/s$.
\item Lại có: $\Delta \ell_0 = \dfrac{mg}{k} = \dfrac{0,2.10}{80} = 0,025\ m = 2,5\ cm$.
\item Đưa vật dọc theo trục của lò xo tới vị trí lò xo nén 1,5 cm rồi buông nhẹ cho vật dao động $\Rightarrow A = 2,5 + 1,5 = 4$ cm
\item Trong 1 chu kì, lo xo bị nén khi đi từ vị trí $M_1$ qua biên âm, tới vị trí $M_2$ như hình vẽ\\
$\Delta \varphi = 2arc\cos(\dfrac{2,5}{4}) = 1,79\ rad$.\\
$\Rightarrow t = \dfrac{\Delta \varphi}{\omega} = \dfrac{1,79}{20} = 0,0896\ s$}
{\begin{tikzpicture}[scale=1,thick,>=stealth']
\tkzInit[ymin=-3,ymax=3,xmin=-3,xmax=3]\tkzClip
\tkzDefPoints{-2.5/0/m, 2.5/0/n, 0/0/o, 0/2.5/p, 0/-2.5/q}
\tkzDefShiftPoint[o](129:2){1}
\tkzDefShiftPoint[o](-129:2){2}
\tkzDefPointBy[projection=onto m--n](1) \tkzGetPoint{h}
\draw(o)circle(2cm);
\tkzDrawSegments[dashed](1,2)
\tkzDrawSegments(p,q)
\tkzDrawSegments[->](m,n)
\tkzDrawSegments[blue,->](o,1 o,2)
\tkzDrawPoints[size=3](o,1,h,2)
\tkzMarkAngles[size=0.7cm,arrows=->](1,o,2)
\node at (1)[above left]{$M_1$};
\node at (2)[below left]{$M_2$};
\node at (2,0)[below right]{$4$};
\node at (h)[below]{$-2,5$};
\end{tikzpicture}}
\end{itemize}
}
%\haidong{9}
\end{vidu} %\dotlineEX{6}
%\loai{Cho thời gian lò xo nén}

%\loai{Tính tỉ số thời gian lò xo dãn - nén}

%\loai{Cho tỉ số thời gian lò xo dãn - nén}
\begin{vidu}%[3P1-10.2K][Loai1]
Con lắc lò xo treo thẳng đứng ở vị trí cân bằng lò xo dãn 5 cm. Lấy $g = 10\ m/s^2$. Biết rằng trong một chu kì thời gian lò xo bị nén bằng một nửa thời gian lò xo dãn. Tốc độ của vật khi đi qua vị trí lò xo không biến dạng là
\choice
{$\dfrac{\sqrt{2}}{2}$ m/s}
{$\dfrac{\sqrt{3}}{2}$ m/s}
{$\dfrac{\sqrt{3}}{3}$ m/s}
{\True $\dfrac{\sqrt{6}}{2}$ m/s}
\loigiai{
\begin{itemize}
\item $\omega  =  10\sqrt{2} $ rad/s.
\item Trong một chu kì, khoảng thời gian lò xo nén là $ \dfrac{T}{3}  \Rightarrow \Delta \ell_0 =   \dfrac{A}{2}  \Rightarrow A = 10\  cm$.
\item Khi lò xo không biến dạng vật có: $|x| = \Delta \ell_0 =   \dfrac{A}{2}  \Rightarrow |v| = \dfrac{\omega A\sqrt{3}}{2} = \dfrac{\sqrt{6}}{2}\ m/s$.
\item 
\item 
\end{itemize}
}
%\haidong{7}
\end{vidu} %\dotlineEX{6} 
%\loai{Tính thời gian lực đàn hồi cùng chiều, ngược chiều với lực kéo về}
\begin{vidu}%[3P1-10.2T][Loai1] 
Một con lắc lò xo treo vào một điểm cố định, dao động điều hòa theo phương thẳng đứng với chu kì 1,8 s. Trong một chu kì, nếu tỉ số của thời gian lò xo dãn với thời gian lò xo nén bằng 2 thì thời gian mà lực đàn hồi ngược chiều lực kéo về là
\choice
{0,5 s}
{0,4 s}
{0,1 s}
{\True 0,3 s}
\loigiai{
\immini{\begin{itemize}
\item Trong 1 chu kì có lúc lò xo dãn, có lúc lò xo nén nên $\Delta \ell_0 < A$. 
\item Theo đề bài, ta có: $\dfrac{\Delta t_{gian}}{\Delta t_{nen}} = \dfrac{\Delta \varphi_{gian}}{\Delta \varphi_{nen}} = 2$.
\item Mà $\Delta \varphi_{gian} +\Delta \varphi_{nen} = 2\pi \Rightarrow \Delta \varphi_{nen} = \dfrac{2\pi}{3} \Rightarrow \Delta \ell_0 = \dfrac{A}{2}$.
\item Trong 1 chu kì, lực đàn hồi ngược chiều lực kéo về khi con lắc đi từ $M_1$ đến $M_2$ và từ $M_3$ đến $M_4$\\
$\Rightarrow \Delta \varphi = \dfrac{\pi}{6} + \dfrac{\pi}{6} = \dfrac{\pi}{3}$\\
$\Rightarrow t = \dfrac{T}{6} = \dfrac{1,8}{6} = 0,3\ s$.
\end{itemize}}{\begin{tikzpicture}[scale=1,thick,>=stealth']
\tkzInit[ymin=-3,ymax=3,xmin=-3,xmax=3]\tkzClip
\tkzDefPoints{-2.5/0/m, 2.5/0/n, 0/0/o, 0/2.5/p, 0/-2.5/q}
\tkzDefShiftPoint[o](90:2){1}
\tkzDefShiftPoint[o](120:2){2}
\tkzDefShiftPoint[o](-120:2){3}
\tkzDefShiftPoint[o](-90:2){4}
\tkzDefPointBy[projection=onto m--n](2) \tkzGetPoint{h}
\draw(o)circle(2cm);
\tkzDrawSegments[dashed](2,3)
\tkzDrawSegments(p,q)
\tkzDrawSegments[->](m,n)
\tkzDrawSegments[blue,->](o,1 o,2 o,4 o,3)
\tkzDrawPoints[size=3](o,1,h,4,2,3)
\tkzMarkAngles[size=0.7cm,arrows=->](1,o,2)
\tkzMarkAngles[size=0.7cm,arrows=->](3,o,4)
\node at (1)[above right]{$M_1$};
\node at (2)[above left]{$M_2$};
\node at (3)[below left]{$M_3$};
\node at (4)[below right]{$M_4$};
\node at (2,0)[above right]{$A$};
\node at (h)[below]{$-\dfrac{A}{2}$};
\end{tikzpicture}}
}
%\haidong{11}
\end{vidu} %\dotlineEX{7}
